\section*{الفصل \ref{ch:g6:symmetry} --- التماثل المحوري}

\begin{solution}{exo:g6:symmetry:1}
النقطة $A'$ على بُعد مربعين إلى \emph{يمين} $d$، في السطر نفسه الذي فيه
$A$؛ والنقطة $B'$ على بُعد خمسة مربعات إلى اليمين، ومربع واحد أعلى --- أي مواضع
مرآتية.
\end{solution}

\begin{solution}{exo:g6:symmetry:2}
العمودي على $d$ المارّ بالنقطة $M$، وقدمه $H$؛ ثم $M'$ في الجهة
الأخرى حيث $HM' = MH = 3$ سم. والمسافة $MM'$ هي
$3 + 3 = 6$ سم.
\end{solution}

\begin{solution}{exo:g6:symmetry:3}
المحور العمودي (مثل A): A, H, I, M, O, T, U, V, W, X, Y.
والمحور الأفقي (مثل B): B, C, D, E, H, I, K, O, X.
والاثنان معًا: H, I, O, X.
(وتتوقف الأجوبة المضبوطة قليلًا على طريقة رسم الحروف.)
\end{solution}

\begin{solution}{exo:g6:symmetry:4}
\omterm{def:g6:shapes:triangles}{المثلث المتساوي الأضلاع}: $3$ محاور. والمستطيل: $2$ (وهما الواسطان
العموديان لضلعيه). \omterm{def:g6:shapes:quadrilaterals}{والمعيّن}: $2$ (قطراه). والمربع: $4$.
والدائرة: بلا نهاية (كل \omterm{def:g6:lines:objects}{مستقيم} يمرّ بالمركز).
\end{solution}

\begin{solution}{exo:g6:symmetry:5}
\emph{خطأ.} فطيّ مستطيل (غير مربع) على طول \omterm{def:g4:geometry:circle}{قطر} لا
يجعل النصفين ينطبقان: فللمثلثين الشكل نفسه
لكنهما يقعان بشكل مختلف --- جرّب ذلك بورقة. والمربع وحده
يعمل فيه القطران محورين.
\end{solution}

\begin{solution}{exo:g6:symmetry:6}
الانعكاسات تحفظ الأطوال والزوايا والمحيطات
(\cref{prop:g6:symmetry:preserved}): إذن $A'B' = 4$ سم،
و $\widehat{B'A'C'} = 70^\circ$، \omterm{def:g4:measure:perimeter}{ومحيط} $A'B'C' = 12$ سم.
\end{solution}

\begin{solution}{exo:g6:symmetry:7}
إنشاء بالفرجار كما في \cref{met:g6:symmetry:bisector}. ومن أجل أيّ
نقطة $M$ على \omterm{def:g6:lines:objects}{المستقيم}، تؤكد المسطرة أن $MA = MB$
(\cref{prop:g6:symmetry:bisector}).
\end{solution}

\begin{solution}{exo:g6:symmetry:8}
كل النقط التي على بُعد متساوٍ من $A$ و $B$ تكوّن
\omterm{def:g6:symmetry:bisector}{الواسط العمودي} \omterm{def:g6:lines:objects}{للقطعة} $[AB]$ (\cref{prop:g6:symmetry:bisector}).
فيمكن أن تُغرس النافورة في أيّ موضع من ذلك \omterm{def:g6:lines:objects}{المستقيم} (داخل الحديقة).
\end{solution}

\begin{solution}{exo:g6:symmetry:9}
كل رأس ينعكس عموديًّا على المحور: فمن أجل المحور
في $45^\circ$، النقطة التي تبعد $k$ مربعًا إلى يمين المحور تسقط
$k$ مربعًا فوقه (والعكس بالعكس) --- فيتبادل الانتقالان الأفقي والعمودي.
اِعكس \omterm{def:g5:solids:def}{الرؤوس} واحدًا واحدًا، ثم
صِلها بالترتيب.
\end{solution}

\begin{solution}{exo:g6:symmetry:10}
\omterm{def:g6:symmetry:def}{مماثلة} $M$ هي التقاطع الثاني للدائرة مع
العمودي على $d$ المارّ بالنقطة $M$ --- وهي تقع على الدائرة. والسبب:
أن الانعكاس يحفظ المسافات ويثبّت $O$ (الواقعة على $d$)، إذن
صورة أيّ نقطة تبعد $r$ عن $O$ تبعد أيضًا
$r$ عن $O$: فالدائرة ترتسم على نفسها.
\end{solution}

\begin{solution}{exo:g6:symmetry:11}
كل نقطة تقاطع $P$ رُسمت بفتحة الفرجار نفسها $\rho$
من $A$ ومن $B$: إذن $PA = \rho = PB$. وحسب
\cref{prop:g6:symmetry:bisector}، النقطة المتساوية البُعد عن $A$ و $B$
تقع على \omterm{def:g6:symmetry:bisector}{الواسط العمودي} \omterm{def:g6:lines:objects}{للقطعة} $[AB]$؛ ونقطتان كهاتين تعيّنان
\omterm{def:g6:lines:objects}{المستقيم}.
\end{solution}

\begin{solution}{exo:g6:symmetry:12}
من أجل أيّ نقطة اِرتداد $P$ على الجدار، $PB = PB'$ (فالانعكاس بالنسبة
إلى $d$ يحفظ المسافات)، إذن طول المسار
$AP + PB = AP + PB'$. والخط المنكسر $A \to P \to B'$ أقصر ما يكون
حين يكون مستقيمًا، أي\ حين تقع $P$ على \omterm{def:g6:lines:objects}{القطعة} $[AB']$. إذن
أفضل نقطة اِرتداد هي تقاطع $[AB']$ مع الجدار.
\end{solution}

\begin{solution}{pb:g6:symmetry:1}
\textbf{1.} إنشاء كما في \cref{met:g6:symmetry:bisector}؛
ومن أجل أيّ نقطة $M$ على \omterm{def:g6:lines:objects}{المستقيم}، تؤكد المسطرة أن $MA = MB$
(\cref{prop:g6:symmetry:bisector}).

\textbf{2.} يحتاج سؤال النافورة إلى الاتجاه الثاني:
\emph{كل} النقط المتساوية البُعد تقع على الواسط، إذن
المواضع الممكنة هي ذلك \omterm{def:g6:lines:objects}{المستقيم} بالضبط ولا شيء غيره. والنقطة
التي ليست على الواسط ليست إذن متساوية البُعد:
فهي أقرب تمامًا إلى إحدى الشجرتين.

\textbf{3.} يشترك المساران في الضلع $[MP]$، و $PB = PA$
لأن $P$ على الواسط. إذن للمسار $M \to P \to B$
الطول نفسه الذي للمسار $M \to P \to A$، وهو $MP + PA$.
والآن $MB$ هو الطريق \omterm{def:g6:lines:objects}{المستقيم} إلى $B$ --- وهو هنا \emph{نفسه}
المسار عبر $P$، إذن $MB = MP + PA$؛ والطريق \omterm{def:g6:lines:objects}{المستقيم} إلى
$A$ أقصر من الالتفاف عبر $P$: أي $MA < MP + PA = MB$.
فكل نقطة في جهة $A$ من الواسط أقرب إلى $A$.

\textbf{4.} الحدّ هو \omterm{def:g6:symmetry:bisector}{الواسط العمودي} \omterm{def:g6:lines:objects}{للقطعة} $[AB]$:
فالمدرسة $A$ تخدم بالضبط نصف المدينة الواقع في جهتها من
ذلك \omterm{def:g6:lines:objects}{المستقيم} (السؤال 3)، والمدرسة $B$ تخدم النصف الآخر، والأطفال
الساكنون على \omterm{def:g6:lines:objects}{المستقيم} نفسه يختارون أيًّا منهما.

\textbf{5.} تقع $O$ على واسط $[AB]$، إذن $OA = OB$؛
وعلى واسط $[BC]$، إذن $OB = OC$
(\cref{prop:g6:symmetry:bisector}).

\textbf{6.} من $OA = OB$ و $OB = OC$ ينتج $OA = OC$. وحسب
الاتجاه الثاني في \cref{prop:g6:symmetry:bisector}، النقطة
المتساوية البُعد عن $A$ و $C$ تقع على \omterm{def:g6:symmetry:bisector}{الواسط العمودي}
\omterm{def:g6:lines:objects}{للقطعة} $[AC]$: فالواسط الثالث يمرّ بالنقطة $O$ دون أن
يُطلب. إذن الوسائط العمودية الثلاثة لأيّ مثلث
\emph{متلاقية}.

\textbf{7.} الدائرة التي مركزها $O$ ونصف قطرها $OA$ تتألف من
كل النقط التي تبعد $OA$ عن $O$
(\cref{def:g6:lines:circle})؛ وبما أن $OB = OA$ و $OC = OA$، فالنقطتان
$B$ و $C$ عليها. مركز واحد، \omterm{def:g3:shapes:shapes}{ونصف قطر} واحد، وثلاث
نقط تُخدَم: فهي الدائرة المارّة بالنقط $A$ و $B$ و $C$.

\textbf{8.} يكفي واسطان لأن الثالث تلقائي
(السؤال 6): فنقطة تقاطعهما تحقّق أصلًا التساويات الثلاثة
بين المسافات. وتُرسم الدائرة بمركز $O$ وبفتح
الفرجار من $O$ إلى أيّ رأس من \omterm{def:g5:solids:def}{الرؤوس} الثلاثة.

\textbf{9.} الوصفة: علّم ثلاث نقط $A$ و $B$ و $C$ على
القوس الباقي؛ واُرسم الواسطين العموديين للوترين
$[AB]$ و $[BC]$؛ فنقطة تقاطعهما هي مركز
الصحن، ومسافتها إلى $A$ هي \omterm{def:g3:shapes:shapes}{نصف القطر}. وهي تنجح لأن
حافة الصحن دائرة تمرّ بالنقط الثلاث المعلَّمة ---
والسؤال 8 يجد \emph{الدائرة} المارّة بها. والقوس
المختبَر يُعاد بناؤه بالضبط.

\textbf{10.} إذا كانت $A$ و $B$ و $C$ مستقيمية، فواسط
$[AB]$ وواسط $[BC]$ كلاهما عمودي على
\omterm{def:g6:lines:objects}{المستقيم} $(AB) = (BC)$ نفسه --- إذن هما متوازيان
(\cref{prop:g6:lines:facts})، وبما أنهما متمايزان فلا يتقاطعان أبدًا.
ولا نقطة على بُعد متساوٍ من الثلاث، ولا دائرة تمرّ
بثلاث نقط مستقيمية.

\textbf{11.} \omterm{def:g4:geometry:polygon}{الخماسي} المنتظم: $5$ محاور، كل واحد يمرّ برأس
ومنتصف الضلع المقابل (\omterm{def:g3:numbers:evenodd}{العدد الفردي}: فكل المحاور
من هذا النوع الواحد). \omterm{def:g4:geometry:polygon}{والسداسي} المنتظم: $6$ محاور --- ثلاثة تمرّ
برأسين متقابلين وثلاثة بمنتصفَي ضلعين متقابلين
(\omterm{def:g3:numbers:evenodd}{العدد الزوجي}: نوعان، نصف ونصف). والنمط: \omterm{def:g4:geometry:polygon}{المضلع}
المنتظم ذو $n$ ضلعًا له $n$ محورًا.

\textbf{12.} القصّ يمرّ بالطبقتين معًا، إذن
تفقد الطبقتان \omterm{def:g6:lines:objects}{قطعتين} متطابقتين بالضبط؛ والنشر هو
الانعكاس بالنسبة إلى خط الطيّ، وهو يرسل حافة كل طبقة
على حافة الأخرى. فالثقب صورته المرآتية نفسها: أي إن خط
الطيّ محور تماثله --- مهما كان المقصوص.

\textbf{13.} للثقب المنشور $2$ محورا تماثل على
الأقلّ، وهما خطّا الطيّ، وتظهر القصّة في $4$ نسخ
(فكل نشر يضاعف الصورة: $1 \to 2 \to 4$). والطيّ
مرات أكثر على نمط قصاصات الثلج يضاعف النسخ والمحاور
أكثر.

\textbf{14.} للشكل $F''$ الاتجاه \emph{نفسه} الذي للشكل $F$ (فقلبتا
المرآة تلغي إحداهما الأخرى)، \omterm{def:g5:solids:volume}{والحجم} نفسه والارتفاع
نفسه --- فهو ببساطة $F$ منزلقًا إلى اليمين. والانزلاق
يقيس $8$ مربعات: أي \emph{ضعف} المسافة بين
المحورين. فالمرآتان المتوازيتان معًا لا تعكسان
البتة: بل \emph{تنقلان}.

\textbf{15.} الشكل $F''$ \omterm{def:g5:solids:volume}{بالحجم} نفسه من جديد، لكنه يظهر الآن
\emph{مقلوبًا}، مُدارًا نصف دورة حول نقطة
تقاطع المحورين: فكل نقطة من $F''$ تقابل قطريًّا
النقطة الموافقة من $F$ عبر ذلك المركز.
فالمرآتان المتعامدتان تتركّبان في \emph{نصف الدورة} ---
وهو التماثل المركزي في \cref{ch:g7:central}.
\end{solution}
